% SEGMENT OLP-0655-S01
% Part: proof-theory
% Chapter: cut-elimination
% Section: ce-largest

\documentclass[../../../include/open-logic-section]{subfiles}

\begin{document}

\olfileid{pt}{cut}{inv}

\olsection{மிக உயர்ந்த தர வெட்டுகளை நீக்குதல்}

\olref[seq][inv]{lem:G3c-invert} இன் நிறுவல் காட்டியபடி, \Log{G3c}
ஓர் அனுமான விதியின் முடிவை !!{prove} முடிந்தால், \RightR{\lexists},
\LeftR{\lforall} ஆகிய இரு விதிகளைத் தவிர, அதன் ஒவ்வொரு
முன்கோளையும் நிறுவலின் !!{height} ஐ உயர்த்தாமலே !!{prove}
முடியும். இதை மேலும் வலுப்படுத்தி, $ \Log{G3c} + \CutCS$ க்கும்
இதே பண்பு பொருந்தும் எனக் காட்டுவோம்.

முன்போலவே, ஒரு \CutCS{} அனுமானத்தின் \emph{வெட்டுத் தரம்}
என்பது அதன் வெட்டு !!{formula} வின் ஆழம்.

\begin{lem}\ollabel{lem:inv-G3c-cut}
$\Log{G3c} + \CutCS$ இல் !!a{proof}~$\pi$ கொண்டு,
\RightR{\lexists}, \LeftR{\lforall} ஆகியவற்றைத் தவிர்த்த
ஓர் அனுமான விதி~$R$ இன் முடிவை !!{prove} முடிந்தால்,
அதன் ஒவ்வொரு முன்கோளையும் !!{proof}~$\pi'$ கொண்டும்
(விதிக்கு இரண்டு முன்கோள்கள் இருந்தால் !!{proof}s~$\pi'$,
$\pi''$ கொண்டும்) !!{prove} முடியும். மேலும்,
(a)~$\pheight{\pi'}\le\pheight{\pi}$
(மற்றும் $\pheight{\pi''}\le\pheight{\pi}$);
(b)~$\pi'$ இல் (மற்றும் $\pi''$ இல்) உள்ள \CutCS{}
அனுமானங்களின் எண்ணிக்கையும் அவற்றின் தரங்களும்~$\pi$ இல்
இருப்பனவே.
\end{lem}

% SEGMENT OLP-0655-S02
\begin{proof}
\cref{pt:seq:inv:lem:G3c-invert,pt:seq:inv:lem:invert-quant}
ஆகியவற்றின் நிறுவல்களை ஆராய்ந்தால் போதும். அவற்றில்
$\pheight{\pi}$ மீது தொகுத்தறிதல் பயன்படுத்தப்பட்டது.
அடிப்படை நேர்வில் ஓர் அடிக்கோளை மற்றோர் அடிக்கோளால்
மாற்றினோம். எனவே (a), (b) உடனே நிறைவேறுகின்றன.
~$\pi$ இன் கடைசி விதி~$R$ ஆக இருந்தால், அதன் உடனடி
உள்-!!{proof}s ஆன~$\pi_i$ களே நமக்கு வேண்டிய
!!{proof}s; இங்கும் (a), (b) நிறைவேறுகின்றன.
மற்ற எல்லா நேர்வுகளிலும், ~$\pi$ இன் உடனடி
உள்-!!{proof}s களுக்கு, அதாவது கடைசி அனுமானம்~$R$
இன் முன்கோள்களுக்கு இட்டுச் செல்பவற்றிற்குத்,
தொகுத்தறிதல் கருதுகோளைப் பயன்படுத்தினோம்; பின்னர்
அதன் விளைவுகளுக்கு அதே விதி $R$ ஐப் பயன்படுத்தினோம்.
புதிய நிறுவலின் உடனடி உள்-!!{proof}s களில் (a), (b)
நிறைவேறுவதால் முழு நிறுவலிலும் அவை நிறைவேறுகின்றன.
கடைசி அனுமானம்~\CutCS{} ஆகும் நேர்வை மட்டும்
இனி உறுதிசெய்ய வேண்டும். இதற்கும் ஓர் எடுத்துக்காட்டையே
பார்ப்போம்: $\pi$ என்பது \LeftR{\land} இன் முடிவான
$!B \land !C, \Gamma \Sequent \Delta$ இன் !!a{proof}
என்றும் அது~\CutCS{} இல் முடிகிறது என்றும் கொள்வோம்:
\[
\AxiomC{}
\RightLabel{$\pi_1$}
\Deduce$!B \land !C, \Gamma \fCenter \Delta, !A$
\AxiomC{}
\RightLabel{$\pi_2$}
\Deduce$!A, !B \land !C, \Gamma \fCenter \Delta$
\RightLabel{\CutCS}
\BinaryInf$!B \land !C, \Gamma \fCenter \Delta$
\DisplayProof
\]
$\pi_1$, $\pi_2$ இரண்டும் $\LeftR{\land}$ இன்
முடிவுகளின் நேர்வுகளுக்கான !!{proof}s என்பதால்,
தொகுத்தறிதல் கருதுகோளை அவற்றிற்குப் பயன்படுத்தலாம்.
அதன் மூலம் முறையே
$!B, !C, \Gamma \fCenter \Delta, !A$,
$!A, !B, !C, \Gamma \fCenter \Delta$
ஆகியவற்றின் !!{proof}s~$\pi_1'$,~$\pi_2'$
கிடைக்கும். \CutCS{} கொண்டு இவற்றை இணைத்து~$\pi'$
ஐப் பெறுகிறோம்:
\[
\AxiomC{}
\RightLabel{$\pi_1'$}
\Deduce$!B, !C, \Gamma \fCenter \Delta, !A$
\AxiomC{}
\RightLabel{$\pi_2'$}
\Deduce$!A, !B, !C, \Gamma \fCenter \Delta$
\RightLabel{\CutCS}
\BinaryInf$!B, !C, \Gamma \fCenter \Delta$
\DisplayProof
\]
(a), (b) இங்கும் நிறைவேறுவது தெளிவு.
\end{proof}

\CutCS{} க்குப் பதிலாக \Cut{} ஐப் பயன்படுத்தினால்
இது செயல்படாது. அப்போது கடைசி !!{proof} இன் முடிவில்,
முன்பக்கத்தில் $!B, !C, \Gamma$ இருமுறையும்
பின்பக்கத்தில் $\Delta$ இருமுறையும் வரும். சுருக்கலின்
ஏற்கத்தக்கதன்மையை நாட வேண்டியிருக்கும்; ஆனால்
\olref[seq][inv]{prop:G3c-cont-adm} இன் நிறுவல்
மீளாக்கத் துணைத்தேற்றத்தைப் பயன்படுத்துகிறது.
அப்படிச் செய்தால் வட்டம்வாதமாகிவிடும்.

% SEGMENT OLP-0655-S03
\cref{pt:cut:inv:lem:inv-G3c-cut} ஐப் பயன்படுத்தி
வெட்டு நீக்குதலுக்கு மாற்று நிறுவல் தரலாம். இதில்
உச்சியிலுள்ள \CutCS{} அனுமானங்களை ஒவ்வொன்றாக
நீக்குவதற்குப் பதிலாக, மிக உயர்ந்த தரமுடைய
\CutCS{} அனுமானங்களை நீக்குகிறோம். அதாவது
$\Log{G3c} + \CutCS$ இல் ஓர் !!a{proof}~$\pi$
கொண்டு, அதில் எங்குள்ள வெட்டாக இருந்தாலும்
அதை நீக்குகிறோம்; சிறிய தரமுடைய வெட்டுகள் அதன்
இடத்தில் வரலாம். வெட்டு எதுவும் மீதமிராதவரை
இதைத் தொடர்கிறோம். நாம் கையாளும் வெட்டுக்கு
மேலே அதே தரமோ அதைவிட உயர்ந்த தரமோ உடைய
வெட்டுகள் இருக்கக்கூடாது என்பதே நிபந்தனை.

\begin{lem}\ollabel{lem:max-cut-red-G3c}
$\Log{G3c}+\CutCS$ இல் உள்ள !!a{proof}~$\pi$,
தரம்~$n$ உடைய~\CutCS{} அனுமானத்தில் முடிந்து,
மற்ற இடங்களில்~$\Log{G3c}$ இன் விதிகளையோ
தரம் $<n$ உடைய \CutCS{} அனுமானங்களையோ
மட்டுமே பயன்படுத்தினால், அதே இறுதித்
தொடரணிக்கான~$\Log{G3c} + \CutCS$ இல்
ஓர் நிறுவல்~$\pi'$ உண்டு; அதிலுள்ள ஒவ்வொரு
\CutCS{} அனுமானத்தின் தரமும்~$<n$.
\end{lem}

\begin{proof}
!!{proof}~$\pi$ இவ்வாறு முடிகிறது:
\[
\AxiomC{}
\RightLabel{$\pi_1$}
\Deduce$\Gamma \fCenter \Delta, !A$
\AxiomC{}
\RightLabel{$\pi_2$}
\Deduce$!A, \Gamma \fCenter \Delta$
\RightLabel{\CutCS}
\BinaryInf$\Gamma \fCenter \Delta$
\DisplayProof
\]
~$!A$ இன் வடிவத்தைப் பொறுத்து நேர்வுகளைப் பிரிப்போம்.

முதலில் $!A$ அணுவாய்பாடு என்க. ~$\pi_1$,
~$\pi_2$ இல் உள்ள எல்லா \CutCS{} அனுமானங்களின்
தரமும் $< \depth{!A} = 0$ ஆக வேண்டும் என்ற
நிபந்தனையால், ~$\pi_1$ அல்லது~$\pi_2$ இல்
\CutCS{} அனுமானம் எதுவும் இருக்க முடியாது.
$ \pheight{\pi_2}$ மீது
தொகுத்தறிந்து $\Gamma \Sequent \Delta$ க்கு
வெட்டற்ற நிறுவல் உண்டு எனக் காட்டுவோம்.
$\pheight{\pi_2} = 0$ எனில்,
$!A, \Gamma \Sequent \Delta$ ஓர் அடிக்கோள்.
$!A$ முதன்மை வாய்பாடல்ல எனில், ஓர் அணு
$!C$ ஆனது~$\Gamma$, $\Delta$ இரண்டிலும்
!!a{element} ஆகும்; ஆகவே
$\Gamma \Sequent \Delta$ தானும் அடிக்கோள்.
$!A$ முதன்மை வாய்பாடு எனில்,
$\Delta = \Delta', !A$. இந்நேர்வில்
$\pi_1$ ஆனது
$\Gamma \Sequent \Delta', !A, !A$
இல் முடிகிறது.
\olref[seq][inv]{prop:G3c-cont-adm}
(சுருக்கலின் ஏற்கத்தக்கதன்மை) பயன்படுத்தி
$\Gamma \Sequent \Delta', !A$ க்கு
வெட்டற்ற நிறுவல் பெறுகிறோம்.
$\pheight{\pi_2} > 0$ எனில் அது
ஒரு விதி~$R$ இல் முடிகிறது. அந்த விதியின்
ஒரு முன்கோள் $!A, \Gamma', \Pi \Sequent
\Delta', \Lambda$ என்ற வடிவில் இருக்கும்;
இங்கு $\Pi$, $\Lambda$ ஆகியவை~$R$ இன்
செயல்படு !!{formula}s.
\olref[seq][adm]{prop:weak-G3c-adm} ஐ
~$\pi_1$, $\pi_2$ ஆகியவற்றிற்குப்
பயன்படுத்தினால்
$\Gamma, \Gamma', \Pi \Sequent \Delta, \Delta', \Lambda, !A$
மற்றும்
$!A, \Gamma, \Gamma', \Pi \Sequent \Delta, \Delta', \Lambda$
ஆகியவற்றின் !!{proof}s கிடைக்கும்.
இவற்றிற்கு தொகுத்தறிதல் கருதுகோள் பொருந்தி,
~$R$ இன் ஒவ்வொரு முன்கோளுக்கும்
$\Gamma, \Gamma', \Pi \Sequent \Delta, \Delta', \Lambda$
இன் !!a{proof} ஐத் தருகிறது.
~$R$ ஐப் பயன்படுத்தினால்
$\Gamma, \Gamma \Sequent \Delta, \Delta$
கிடைக்கும்; பின்னர்
\olref[seq][inv]{prop:G3c-cont-adm}
மூலம் $\Gamma \Sequent \Delta$ க்கான
வெட்டற்ற நிறுவல் கிடைக்கும்.

% SEGMENT OLP-0655-S04
$!A$ அணுவாய்பாடல்ல எனில், ~$!A$ இன்
!!{main operator} ஐப் பொறுத்து நேர்வுகளைப்
பிரிப்போம். ஒவ்வொரு நேர்விலும்,
$!A$ முதன்மை வாய்பாடாக உள்ள இடப்புற,
வலப்புற விதிகளின் முன்கோள்களுக்கான
!!{proof}s பெற \olref{lem:inv-G3c-cut}
ஐப் பயன்படுத்துகிறோம். பின்னர்
\olref[top]{lem:cut-adm-G3c} இன்
நிறுவலில் உள்ள E நேர்வைப் போலத் தொடர்கிறோம்.
\begin{enumerate}
  \item $!A \ident \lnot !B$.
  !!{proof}s~$\pi_1$, $\pi_2$ முறையே
  $\Gamma \Sequent \Delta, \lnot !B$,
  $\lnot !B, \Gamma \Sequent \Delta$
  ஆகியவற்றில் முடிகின்றன.
  \olref{lem:inv-G3c-cut} மூலம்
  $!B, \Gamma \Sequent \Delta$,
  $\Gamma \Sequent \Delta, !B$
  ஆகியவற்றின் !!{proof}s~$\pi_1'$, $\pi_2'$
  கிடைக்கின்றன. !!{proof}~$\pi'$ இதுவே:
  \[
  \AxiomC{}
  \RightLabel{$\pi_2'$}
  \Deduce$\Gamma \fCenter \Delta, !B$
  \AxiomC{}
  \RightLabel{$\pi_1'$}
  \Deduce$!B, \Gamma \fCenter \Delta$
  \RightLabel{\CutCS}
  \BinaryInf$\Gamma \fCenter \Delta$
  \DisplayProof
  \]

% SEGMENT OLP-0655-S05
  \item $!A \ident (!B \lor !C)$.
  !!{proof}s~$\pi_1$, $\pi_2$ முறையே
  $\Gamma \Sequent \Delta, !B \lor !C$,
  $!B \lor !C, \Gamma \Sequent \Delta$
  ஆகியவற்றில் முடிகின்றன.
  ~$\pi_1$ க்கு \RightR{\lor} இன்
  மீளாக்கத்தை வழங்கும்
  \olref{lem:inv-G3c-cut} ஐப் பயன்படுத்தி
  $\Gamma \Sequent \Delta, !B, !C$
  இன் !!a{proof}~$\pi_1'$ பெறுகிறோம்.
  ~$\pi_2$ க்கு \LeftR{\lor} இன்
  மீளாக்கத்தை வழங்கும்
  \olref{lem:inv-G3c-cut} மூலம்
  $!B, \Gamma \Sequent \Delta$ இன்
  !!a{proof}~$\pi_2'$ உம்
  $!C, \Gamma \Sequent \Delta$ இன்
  !!a{proof}~$\pi_2''$ உம் கிடைக்கின்றன.
  $\pi_2''$ க்கு
  \olref[seq][adm]{prop:weak-G3c-adm}
  ஐப் பயன்படுத்தி
  $!C, \Gamma \Sequent \Delta, !B$
  இன் !!a{proof}~$\pi_2'''$ உம்
  பெறுகிறோம். !!{proof}~$\pi'$ இதுவே:
  \[
  \AxiomC{}
  \RightLabel{$\pi_1'$}
  \Deduce$\Gamma \fCenter \Delta, !B, !C$
  \AxiomC{}
  \RightLabel{$\pi_2'''$}
  \Deduce$!C, \Gamma \fCenter \Delta, !B$
  \RightLabel{\CutCS}
  \BinaryInf$\Gamma \fCenter \Delta, !B$
  \AxiomC{}
  \RightLabel{$\pi_2'$}
  \Deduce$!B, \Gamma \fCenter \Delta$
  \RightLabel{\CutCS}
  \BinaryInf$\Gamma \fCenter \Delta$
  \DisplayProof
  \]

% SEGMENT OLP-0655-S06
  \item $!A \ident (!B \land !C)$.
  பயிற்சி.
  \item $!A \ident (!B \lif !C)$.
  !!{proof}s~$\pi_1$, $\pi_2$ முறையே
  $\Gamma \Sequent \Delta, !B \lif !C$,
  $!B \lif !C, \Gamma \Sequent \Delta$
  ஆகியவற்றில் முடிகின்றன.
  ~$\pi_1$ க்கு \RightR{\lif} இன்
  மீளாக்கத்தை வழங்கும்
  \olref{lem:inv-G3c-cut} மூலம்
  $!B, \Gamma \Sequent \Delta, !C$
  இன் !!a{proof}~$\pi_1'$ கிடைக்கிறது.
  ~$\pi_2$ க்கு \LeftR{\lif} இன்
  மீளாக்கத்தை வழங்கும்
  \olref{lem:inv-G3c-cut} மூலம்
  $\Gamma \Sequent \Delta, !B$
  இன் !!a{proof}~$\pi_2'$ உம்
  $!C, \Gamma \Sequent \Delta$
  இன் !!a{proof}~$\pi_2''$ உம்
  கிடைக்கின்றன.
  $\pi_2''$ க்கு
  \olref[seq][adm]{prop:weak-G3c-adm}
  ஐப் பயன்படுத்தி
  $!C, !B, \Gamma \Sequent \Delta$
  இன் !!a{proof}~$\pi_2'''$ உம்
  பெறுகிறோம். !!{proof}~$\pi'$ இதுவே:
  \[
  \AxiomC{}
  \RightLabel{$\pi_2'$}
  \Deduce$\Gamma \fCenter \Delta, !B$
  \AxiomC{}
  \RightLabel{$\pi_1'$}
  \Deduce$!B, \Gamma \fCenter \Delta, !C$
  \AxiomC{}
  \RightLabel{$\pi_2'''$}
  \Deduce$!C, !B, \Gamma \fCenter \Delta$
  \RightLabel{\CutCS}
  \BinaryInf$!B, \Gamma \fCenter \Delta$
  \RightLabel{\CutCS}
  \BinaryInf$\Gamma \fCenter \Delta$
  \DisplayProof
  \]

% SEGMENT OLP-0655-S07
  \item $!A \ident \lforall[x][!B(x)]$.
  !!{proof}s~$\pi_1$, $\pi_2$ முறையே
  $\Gamma \Sequent \Delta, \lforall[x][!B(x)]$,
  $\lforall[x][!B(x)], \Gamma \Sequent \Delta$
  ஆகியவற்றில் முடிகின்றன.
  \olref{lem:inv-G3c-cut} மூலம்
  $\Gamma \Sequent \Delta, !B(c)$ இன்
  !!a{proof}~$\pi_1'$ கிடைக்கிறது.

  இப்போது !!{proof}~$\pi_2$ ஐக் கருதுவோம்.
  அது $\lforall[x][!B(x)], \Gamma \Sequent
  \Delta$ இல் முடிகிறது.
  ~$\pi_2$ இன் இறுதித் தொடரணியிலிருந்து
  மேல்நோக்கிச் செல்லும்போது, இடப்புறத்தில்
  வரும் $\lforall[x][!B(x)]$ இன்
  எந்த நிகழ்வு, ~$\pi_2$ இன்
  இறுதித் தொடரணியிலுள்ள
  $\lforall[x][!B(x)]$ இன் நிகழ்வுக்கு
  “இட்டுச் செல்கிறதோ” அந்த ஒவ்வொன்றையும்
  குறியிடுவோம். குறிப்பாக:
  (a)~$\pi_2$ இன் இறுதித் தொடரணியில்
  $\lforall[x][!B(x)]$ ஐக் குறியிடுக.
  (b)~ஒரு விதியின் முடிவின் சூழலில்
  $\lforall[x][!B(x)]$ வந்து
  குறியிடப்பட்டிருந்தால், அதன்
  முன்கோளில் அல்லது முன்கோள்களில்
  ஒத்த நிகழ்வைக் குறியிடுக.
  (c)~ஒரு \LeftR{\lforall} விதியின்
  முடிவில் $\lforall[x][!B(x)]$
  முதன்மை வாய்பாடாக வந்து
  குறியிடப்பட்டிருந்தால், அதன்
  முன்கோளிலுள்ள ஒத்த செயல்படு
  நிகழ்வுகளான
  $\lforall[x][!B(x)]$,
  $!B(t)$ ஆகிய இரண்டையும் குறியிடுக.

  இவ்வாறு குறியிடப்பட்ட
  $\lforall[x][!B(x)]$ இன் எல்லா
  நிகழ்வுகளையும் கருதுவோம்.
  அவற்றுள் எதுவும் \LeftR{\lforall}
  அல்லாத விதியில் செயல்படுவதில்லை;
  \LeftR{\lforall} இன் செயல்படு
  !!{formula}s மட்டுமே குறியிடப்படும்.
  ~$\pi_2$ இல் குறியிடப்பட்ட
  $\lforall[x][!B(x)]$ இன்
  எல்லா நிகழ்வுகளையும் நீக்குக.
  இதனால் சரியான !!{proof} கிடைத்தால்,
  குறியிடப்பட்ட $\lforall[x][!B(x)]$
  நிகழ்வுகள் ஒருபோதும்
  செயல்படவில்லை; அவை அனைத்தும்
  நேரடியாக அடிக்கோள்களிலிருந்து
  வந்தன. அந்த !!{proof} ஆனது
  $\Gamma \Sequent \Delta$ இல்
  முடிவதால்~$\pi$ ஐ அதனால்
  மாற்றலாம். மிக உயர்ந்த தரமுடைய
  ஒரு வெட்டை இவ்வாறு நீக்கியிருக்கிறோம்.

% SEGMENT OLP-0655-S08
  மாறாக, சில \LeftR{\lforall}
  அனுமானங்களில் செயல்பட்ட
  $\lforall[x][!B(x)]$
  !!{formula}s ஐ நாம் நீக்கியதால்,
  கிடைத்தது சரியான நிறுவலாக
  இல்லாமலும் இருக்கலாம். அப்போது
  அவை பின்வரும் வடிவிலான
  “அனுமானங்களாக” மாறியிருக்கும்:
  \[
  \AxiomC{}
  \RightLabel{$\pi_t$}
  \Deduce$!B(t), \Pi \fCenter \Lambda$
  \RightLabel{\LeftR{\lforall}}
  \UnaryInf$\Pi \fCenter \Lambda$
  \DisplayProof
  \]
  இவ்வகையில் உச்சியில்
  முதலில் வரும் ஒவ்வோர்
  “அனுமானத்தையும்” பின்வருவதாக
  மாற்றுக:
  \[
  \AxiomC{}
  \RightLabel{$\Subst{\pi_1'}{t}{c}[\Pi \Sequent \Lambda]$}
  \Deduce$\Gamma, \Pi \fCenter \Delta, \Lambda, !B(t)$
  \AxiomC{}
  \RightLabel{$\pi_t[\Gamma \Sequent \Delta]$}
  \Deduce$!B(t), \Gamma, \Pi \fCenter \Delta, \Lambda$
  \RightLabel{\CutCS}
  \BinaryInf$\Gamma, \Pi \fCenter \Delta, \Lambda$
  \DisplayProof
  \]
  அதன் கீழுள்ள ஒவ்வொரு தொடரணியின்
  இடப்புறத்தில் $\Gamma$ வையும்
  வலப்புறத்தில் $\Delta$ வையும்
  சேர்க்குக. முன்கோளில்
  குறியிடப்பட்ட $!B(t)$ உள்ள
  எல்லா \LeftR{\lforall}
  “அனுமானங்களும்” ஏதேனும்
  $!B(t)$ மீதான \CutCS{}
  அனுமானத்தால் மாற்றப்படும் வரை
  இதைத் தொடர்க.
  இப்போது கிடைக்கும் !!{proof}
  சரியான !!{proof} ஆகும்; அதன் இறுதித்
  தொடரணி
  $\Gamma, \dots, \Gamma \Sequent
  \Delta, \dots, \Delta$ என்ற
  வடிவில் இருக்கும். சுருக்கலின்
  ஏற்கத்தக்கதன்மையைப் பயன்படுத்தி
  இதனை~$\Gamma \Sequent \Delta$
  இன் !!a{proof} ஆக மாற்றி,
  அதனால்~$\pi$ ஐ மாற்றுக.
  $\lforall[x][!B(x)]$ மீதான
  ஒரு வெட்டுக்குப் பதிலாக,
  $!B(t)$ மீதான வெட்டுகளை
  (பலவாகவும் இருக்கலாம்)
  அமைத்துள்ளோம்; அவற்றின்
  தரம் அனைத்தும் குறைவு.
  ஏற்கெனவே $\pi_1$ அல்லது
  $\pi_2$ இல் இருந்த வெட்டுகள்
  தொடர்ந்து இருக்கும்; அவற்றின்
  தரமும் குறைவுதான்.
\end{enumerate}
\end{proof}

பதிப்புக் குறிப்பு: இருத்தல் அளவையுள்ள வெட்டு வாய்பாட்டின் நேர்வு
மூல நிறுவலின் நேர்வாய்வில் எழுதப்படவில்லை. எனவே மேலுள்ள
நேர்வாய்வு மட்டும் துணைத்தேற்றத்தின் எல்லா வாய்பாடுகளுக்கும்
முழு நிறுவலாகாது.

% SEGMENT OLP-0655-S09
\begin{prob}
\olref{lem:max-cut-red-G3c} இன்
நிறுவலில் $!A \ident (!B \land !C)$
என்ற நேர்வை உறுதிசெய்க.
அதாவது, கீழ்க்கண்ட வடிவில்
முடியும் !!a{proof}~$\pi$ இருந்தால்
\[
\AxiomC{}
\RightLabel{$\pi_1$}
\Deduce$\Gamma \fCenter \Delta, !B \land !C$
\AxiomC{}
\RightLabel{$\pi_2$}
\Deduce$!B \land !C, \Gamma \fCenter \Delta$
\RightLabel{\CutCS}
\BinaryInf$\Gamma \fCenter \Delta$
\DisplayProof
\]
மேலும் $\pi_1$, $\pi_2$ இரண்டிலும்
$<\depth{!B \land !C}$ தரமுள்ள
வெட்டுகள் மட்டுமே இருந்தால்,
$\Gamma \Sequent \Delta$ இன்
!!a{proof}~$\pi'$ ஒன்றில்
$<\depth{!B \land !C}$ தரமுள்ள
வெட்டுகள் மட்டுமே இருக்கும்
எனக் காட்டுக.
\end{prob}

% \begin{prob}
% $!A \ident \lexists[x][!B(x)]$ என்ற நேர்வை
% \olref{lem:inv-G3c-cut} இன் நிறுவலில் உறுதிசெய்க.
% \end{prob}

% SEGMENT OLP-0655-S10
\olref{lem:max-cut-red-G3c} மூலம்,
ஓர் !!a{proof} இல் மிக உயர்ந்த
தரமுடைய \CutCS{} அனுமானத்தை
அதற்குக் குறைந்த தரமுள்ள
\CutCS{} அனுமானங்களால்
மாற்ற முடியும். இதே
துணைத்தேற்றத்தை மீண்டும்
பயன்படுத்தி, மிக உயர்ந்த தரமுள்ள
\CutCS{} இல் முடியும் உச்சி
உள்-!!{proof}s களுக்கு அதை
ஒவ்வொன்றாகப் பயன்படுத்துவதன்
மூலம், ஓர் !!a{proof} இலுள்ள
எத்தனை \CutCS{} அனுமானங்களையும்
நீக்க முடியும்.

\begin{thm}\ollabel{thm:cut-elim-G3c-largest}
$\Log{G3c} + \CutCS$ இல் உள்ள
எந்த !!{proof}~$\pi$ ஐயும்
\Log{G3c} இல் ஒரு நிறுவலாக
மாற்றலாம்.
\end{thm}

% SEGMENT OLP-0655-S11
\begin{proof}
  முதலில்~$\pi$ இல் மிக உயர்ந்த
  தரமுள்ள எல்லா வெட்டுகளையும்
  நீக்க முடியும் எனக் காட்டுவோம்.
  ~$\pi$ இல் உள்ள வெட்டுகளின்
  மிக உயர்ந்த தரம்~$d$ எனவும்,
  தரம்~$d$ உள்ள வெட்டுகளின்
  எண்ணிக்கை~$n$ எனவும் கொள்வோம்.
  ~$n$ மீது தொகுத்தறிந்து நிறுவுவோம்.
  $n=0$ எனில் செய்ய எதுவும் இல்லை.
  $n>0$ எனில், தரம்~$d$ உடைய
  உச்சியிலுள்ள ஒரு வெட்டைத்
  தேர்ந்து, அதில் முடியும்
  உள்-!!{proof}~$\pi_1$ ஐக் கொள்வோம்.
  \olref{lem:max-cut-red-G3c}
  மூலம் அதே இறுதித் தொடரணிக்கான
  !!a{proof}~$\pi_1'$ உண்டு;
  அதில் உள்ள எல்லா வெட்டுகளின்
  தரமும் $<d$. ~$\pi$ இல்
  $\pi_1$ க்குப் பதிலாக
  $\pi_1'$ ஐ வைக்கிறோம்.
  விளைவாக, தரம்~$d$ உள்ள
  வெட்டுகள் $n-1$ மட்டுமே
  கொண்ட !!a{proof} கிடைக்கிறது;
  அதற்குத் தொகுத்தறிதல்
  கருதுகோள் பொருந்தும்.

  இப்போது~$d$ மீது தொகுத்தறிந்து
  தேற்றத்தைப் பெறுகிறோம்.
  $d=0$ எனில் எல்லா
  வெட்டுகளும் அணுவாய்பாடுகள்
  மீதானவை; முந்தைய முடிவினால்
  அவை அனைத்தையும் நீக்கலாம்.
  $d>0$ எனில், முந்தைய முடிவு
  தரம் $d$ உள்ள எல்லா
  வெட்டுகளும் தரம்~$<d$
  உள்ள வெட்டுகளால் மாற்றப்பட்ட
  ஒரு நிறுவலைத் தருகிறது.
  ~$\pi$ இன் மிக உயர்ந்த
  வெட்டுத் தரம் $d$ என்பதால்,
  இப்போது மிக உயர்ந்த
  வெட்டுத் தரமே $<d$
  உள்ள ஒரு நிறுவல்
  கிடைத்துள்ளது. அதற்குத்
  தொகுத்தறிதல் கருதுகோள் பொருந்தும்.
\end{proof}

\end{document}
